\section{A compact nine-positive companion}
\label{sec:compact-nine}

The supplied companion polynomial has nine strictly positive unknowns,
exact total degree $1\,834\,516$, and a live arithmetic circuit of $453$
binary additions, subtractions and multiplications.  At the inspected
ProveIt revision, package59's corresponding point is
$(9,1\,910\,404,473)$.  Thus this component saves $75\,888$ degrees and
$20$ operations.  It also binds the source to the complete compact44
coefficient tuple, whose four unsigned binary encodings occupy
$295\,388$ bytes, compared with $533\,099$ for the preceding tuple.
These are distinct measures: the constants remain large, and operation
count is not bit complexity.  This is a local Pareto improvement, not a
claim of optimality or a major advance in universal variable counts.

The authoritative equation is the ordinary integer polynomial represented
by the supplied arithmetic DAG:
\begin{center}
\small\texttt{nine/generated/split\_116\_positive\_global.json}.
\end{center}
Its fixed
leaves are completely materialized and authenticated by
\texttt{nine/fixed/BINDING.json}.  Intermediate register names below are
abbreviations, not quantified variables.  In particular, all displayed
exponents are fixed integers.  The nine inputs, in order, are
\[
 b,g,h,s,w,C,i,j,n\in\mathbb Z_{>0}.
\]
Using the identified inherited native compiler, its compact whole-tuple
adapter, the calibration from K215 to finite refutations and the Jones ratio theorem,
ZFC proves
\[
 (\exists b,g,h,s,w,C,i,j,n>0\;\mathcal F=0)
 \quad\Longleftrightarrow\quad \neg\operatorname{Con}(\mathrm{ZFC}).
\]
The added reduction is ordinary mathematics formalizable in ZFC.  The
package supplies exact finite certificates, not a new end-to-end
Lean/Rocq proof or an exported object-language ZFC derivation.

\subsection{The source specialization and exact target certificate}

Use the $31$-positive-coordinate, $22$-row source in
\texttt{nine/data/parent\_system60.json}.  Its four symbolic coefficient
ports are $\beta,K_0,M_C,M_F$.  The complete compact44 tuple has bit
lengths
\[
 (590069,592884,590069,590070),
\]
with aliases $K_0=K$ and $M_F=M_{F,\mathrm{source}}$; the last port is
already shifted by the native base minus one.  The four payloads are
supplied in \texttt{nine/coefficients/}.  The supplied native builder
reconstructs them exactly from the literal transition data.  The finite
verification includes all $66062$ native basis identities and the
explicit rank controller.

The old package59 source must be adjusted before this tuple is used.
With all four ports symbolic, exact expansion gives
\[
 r_{15}^{\mathrm{new}}=r_{15}^{\mathrm{old}}-\beta J,
 \qquad r_i^{\mathrm{new}}=r_i^{\mathrm{old}}\quad(i\ne15).
\]
This is the justified transport shift $z_{\mathrm{old}}=z+1$; its
completeness uses the inherited source bound $z_{\mathrm{old}}\ge2$ on
zeros.  It is not an arbitrary positive-coordinate substitution.  Apply
the invertible row operations
\[
 \rho_{14}=r_{14}+r_2,\qquad \rho_{15}=r_{15}+r_3,
 \qquad \rho_i=r_i\ (i\notin\{14,15\}).
\]
The resulting source has $89$ records: $7$ constant, $39$ linear and
$43$ quadratic.  In its original variable order choose weights
\[
\begin{split}
(w_i)={}&(1,7,61,36,16,85,20,67,43,11,34,13,96,80,29,50,\\
 &2,109,110,112,116,91,53,55,27,101,92,9,107,37,60).
\end{split}
\]
Assign target exponents in the following paired order:
\[
\begin{split}
\text{rows}:\;&(2,15,8,17,13,18,12,20,21,4,22,9,11,19,1,6,10,16,3,14,7,5),\\
\text{targets}:\;&(28,124,325,423,598,658,775,904,941,955,1027,1236,\\
 &1245,1313,1323,1530,1571,1584,1637,1787,2014,2036).
\end{split}
\]
Write $q_j$ for the target of row $j$, and set
$U_y(T)=\sum_i y_iT^{w_i}$.  For a coefficient $a_{j,\alpha}$ of
$y^\alpha$ in $\rho_j$, let $d=|\alpha|$ and
$e(\alpha)=\sum_i\alpha_iw_i$.  Define
\[
 c_{j,\alpha}=
 \begin{cases}
 2a_{j,\alpha},&d=0,1,\\
 a_{j,\alpha}\prod_i\alpha_i!,&d=2,
 \end{cases}
 \qquad
 D_d(T)=\sum_{j,\,|\alpha|=d}
       c_{j,\alpha}T^{q_j-e(\alpha)}.
\]
The supplied literal sparse kernels and target certificate prove
\[
 [T^{q_j}]\bigl(D_2(T)U_y(T)^2+D_1(T)U_y(T)+D_0(T)\bigr)
       =2\rho_j(y) \quad (1\le j\le22).
\]
Indeed, among all $22544$ same-degree source/expansion interactions,
only the $89$ intended contributions hit selected targets.  Squares
have multiplicity one in $U_y^2$ and cross terms multiplicity two, so
the stated weighting recovers twice the source coefficient.  This is a
formal coefficient identity, with arbitrary witnesses and arbitrary
coefficient ports.  The resulting degrees are
\[
 (\deg D_0,\deg D_1,\deg D_2)=(1571,1934,1836),\quad W=116,
 \quad H=\max_d(\deg D_d+dW)=2068.
\]
The leading coefficient of $D_2$ is exactly $+1$.  There are no negative
kernel exponents.  The maximal target is $2036$, leaving the required
headroom below the convolution height.

\subsection{Source words and the exact binary threshold}

Replacing native-port powers by their strict bit bounds and summing
absolute values gives the sharp common coefficient majorant
$M=2^{1182954}$: every actual weighted coefficient has magnitude strictly
less than $M$.  Set
\[
 \kappa=1182966,\quad a=2^\kappa=4096M,\quad B=ab^2,
 \quad \ell=\sum_iB^{w_i},\quad U_s=\ell+g.
\]
Define the following integer polynomials:
\[
\begin{aligned}
 G_s&=(D_2(B)U_s+D_1(B))U_s+D_0(B),&
 T_*&=\sum_jB^{q_j},\\
 N_0&=B^{117},&Z&=B^{2069},\\
 \mathrm{Mask}&=N_0-1-(b-1)\ell,&
 J_*&=2G_s+T_*,\\
 S_*&=g+N_0J_*,&
 T_{**}&=\mathrm{Mask}+N_0(B-2)T_*.
\end{aligned}
\]
For every $b,g\ge1$, we have $U_s\ge B^{116}$ and
$D_2(B)>B^{1836}/2$.  The two potentially negative contributions, relative
to $B^{1836}U_s^2$, are bounded by
\[
 \frac{|D_1U_s|}{B^{1836}U_s^2}<\frac{2M}{B^{18}}<\frac18,
 \qquad
 \frac{|D_0|}{B^{1836}U_s^2}<\frac{2M}{B^{497}}<\frac18.
\]
Thus $G_s>0$ unconditionally; also $0\le\mathrm{Mask}<N_0$, so the source
words are nonnegative before any finalizer consequence is used.
On the full soundness range $g<4bB^{116}$, one has $U_s<6bB^{116}$ and
\[
 \frac{2G_s}{Z}<\frac{86Mb^2}{B-1}<\frac1{32},\qquad
 T_*<\frac Z{64},\quad(B-2)T_*<\frac Z{64},\quad g<N_0.
\]
Consequently $S_*+T_{**}<B^{2186}$.  For binary $b$, the low mask forces
$g=\sum_i z_iB^{w_i}$ with $0\le z_i<b$, so $U_s$ encodes the positive
witnesses $y_i=z_i+1$.  Every convolution coefficient is bounded in
magnitude by $993Mb^2<B/4$.  The inherited signed-borrow padding argument,
now with the certified target identities and these numerical bounds,
therefore gives
\[
 S_*\mathbin{\&}T_{**}=0
 \quad\Longleftrightarrow\quad
 g\text{ encodes a positive common source zero}.
\]
A positive source zero supplies such a word for every sufficiently large
binary $b$, with $0<g<bB^{116}$.  The strict positivity follows because
the all-one source assignment has $r_5=1$.

Now put
\[
 L=bB^{2186},\quad
 R=L(S_*+T_{**}+1)-S_*-1,\quad
 N=2^{8872234}B^{1155}.
\]
Unconditionally $R,N\ge8$ and $b\le R,N$.  If $b=2^k$, write
$v=1182966+2k$ and $u=k+2186v$, so $B=2^v$ and $L=2^u$.
The bounded source words give the two-block expansion
\[
 R=(S_*+T_{**})L+(L-1-S_*).
\]
With $\operatorname{pop}$ denoting binary population count and
$\tau=\operatorname{pop}(S_*)+\operatorname{pop}(T_{**})
-\operatorname{pop}(S_*+T_{**})$, this implies
\[
 \operatorname{pop}(R)=u+\operatorname{pop}(T_{**})-\tau.
\]
The low mask contributes $117v-31k$ bits and the $22$ high digits
contribute $22v-22$ bits.  Hence the exact calibration is
\[
 u+\operatorname{pop}(T_{**})=2325v-30k-22
 =2(8872234+1155v)=2\log_2N.
\]
Since $\tau\ge0$, with equality exactly for binary-disjoint words,
Legendre's identity
$v_2\binom{2R}{R}=\operatorname{pop}(R)$ proves
\[
 N^2\mid\binom{2R}{R}
 \quad\Longleftrightarrow\quad S_*\mathbin{\&}T_{**}=0.
\]
The case $k=0$ is included and cannot produce a final zero, since then
$\mathrm{Mask}=N_0-1$ and $0<g<N_0$.

\subsection{The explicit final equation}

The remaining formulas are the inherited free-coordinate finalizer,
applied with the parameters just proved.  Write $j_0=j-1\ge0$ and define
\[
\begin{aligned}
 U&=N^2w,&Y&=N^2s,&M_0&=RY,\\
 P&=2M_0^2U,&K&=R+1+h(P-1),&A&=M_0(U+1),\\
 B_1&=2R+1,&W_1&=bw,&J&=4A-5,\\
 Q&=3W_1C-2(W_1^2-1),&&&D&=(A^2-1)C^2+1,\\
 k_0&=iJC,&E&=k_0DC,&F_1&=(A^2-1)E^2+1,\\
 G&=A+F_1(CD+1-A),&&&H_1&=C+F_1(B_1+j_0C),\\
 I&=(G^2-1)H_1^2+1,&a_0&=(P^2-1)K^2+1,&a_1&=DF_1I.
\end{aligned}
\]
Finally set
\[
\begin{aligned}
 t&=bB^{116},&m&=K^2(4t-g)-16t(C-KY)^2,\\
 V_0&=PK,&V_1&=k_0(8D^2-4D-1)GH_1,\\
 \Omega&=Jm,&E_0&=J(n+W_1^2)+Q,\\
 \Theta&=E_0-\Omega(V_0+V_1).
\end{aligned}
\]
The equation is
\begin{equation}
 \boxed{\mathcal F=
 [\Theta^2-\Omega^2(a_0+64a_1)]^2
       -256\Omega^4a_0a_1=0.}
 \label{eq:compact-nine}
\end{equation}
The inherited proof applies because $R,N\ge8$, $b\le R,N$, the source
words are unconditionally nonnegative, and the source completeness bound
$g<bB^{116}$ has the required factor-$32$ margin.
The important positivity step occurs before square-root integrality:
\[
 V_0^2-a_0=K^2-1>0,
 \quad
 V_1^2-64a_1=
 [(8D+1)k_0^2-64D](GH_1)^2+64DF_1(H_1^2-1)>0.
\]
Also $E_0>0$.  A norm zero with $m=0$ would give $E_0^4>0$;
otherwise its real factorization forces $m>0$.  This yields
$|C-KY|<K/2$ and the needed carrier bound.  The rational-root separation
then makes both square roots integral, gives $J\mid Q$, and invokes the
Jones ratio theorem to recover binary $b$, the source range and binomial
divisibility.  Conversely the source and ratio witnesses yield an integer
$n>0$ by the explicit bound $W_1^2+Q/J<2D<V_1$.  Thus no unproved cofinality
claim about $j$ or $n$ enters the construction.  The complete inherited
argument and the new applicability checks are retained in
\texttt{nine/PROOF.md} and its pinned upstream dependencies.

\subsection{Exact degree, cost, and verification}

The source heads give
\[
 \deg S_*=4370,\quad\deg T_{**}=4308,
 \quad r=\deg R=8743,\quad q=\deg N=2310.
\]
Let $\Delta_1=V_1^2-64a_1$.  To resolve the apparent leading cancellation
in the norm, use the exact identity
\[
 \Theta^2-\Omega^2(a_0+64a_1)
 =E_0^2-2E_0\Omega(V_0+V_1)
  +\Omega^2\bigl[(K^2-1)+\Delta_1+2V_0V_1\bigr].
\]
The separate degree certificate has a positive diagonal leading term and
no unresolved positive-degree cancellation.  The diagonal reaches the
multivariate upper bound, proving the exact degree
\[
 104r+400q+8W+316=1\,834\,516.
\]
The initial Horner circuit has $467$ operations.  Its $102$ pure powers of
$B$ are replaced by a shared $88$-gate addition chain.  Every chain row
$e=a+b$ certifies $B^e=B^aB^b$, and exact structural normalization checks
the complete rewritten circuit, including the definition of $B$.  The
final tally is $162$ additions, $22$ subtractions and $269$
multiplications, totaling $453$.  Scalar multiplication and every power
multiplication are included; constant leaves and reuse are free.

Replacing every positive input by a nonnegative input plus one gives the
supplied nine-nonnegative companion, with the same degree and $460$
operations.  No nine-variable unrestricted-integer claim is made.
The $12$ final nonliteral leaves occupy $2\,217\,161$ raw magnitude bytes
or $47\,364$ deterministic gzip bytes; the signs and hashes are explicit.

The command \texttt{python3 nine/reproduce.py} reconstructs the entire
component offline using the Python standard library.  It regenerates the
source and both domains, verifies all target identities, reconstructs the
native tuple, materializes the derived leaves, and checks the power-chain
and exact-degree certificates.  A second finite checker imports no
project builder or verifier.  All authored checks use exceptions and
remain enabled under \texttt{python3 -O}.  Modular and small binary-block
fixtures are separately labeled as finite transcription checks; they do
not establish the unbounded consistency semantics.
